% Generated draft: requires mathematical and source review.
\documentclass[11pt]{article}
\usepackage[margin=25mm]{geometry}
\usepackage{fontspec}
\setmainfont{Latin Modern Roman}
\usepackage{amsmath,amssymb,mathtools,unicode-math}
\setmathfont{Latin Modern Math}
\usepackage{enumitem,needspace,xurl,hyperref}
\setlength{\parindent}{0pt}
\setlength{\parskip}{5pt}
\setlength{\emergencystretch}{3em}
\newcommand{\N}{\mathbb{N}}
\newcommand{\Z}{\mathbb{Z}}
\newcommand{\Q}{\mathbb{Q}}
\newcommand{\R}{\mathbb{R}}
\newcommand{\C}{\mathbb{C}}
\newcommand{\F}{\mathbb{F}}
\newcommand{\A}{\mathbb{A}}
\newcommand{\PP}{\mathbb{P}}
\newcommand{\E}{\mathbb{E}}
\newcommand{\eps}{\varepsilon}
\newcommand{\dd}{\,\mathrm{d}}
\begin{document}
\section*{UnsolvedMath ID: 30004661}
\textbf{Source problem code:} catalogue entry
\subsection{Definitions and mathematical statement}
Let \(k\) be a field, with no restriction on its characteristic. A unipotent algebraic \(k\)-group \(U\) is called \(k\)-wound if every \(k\)-morphism \(\mathbb A_k^1\to U\) is constant. The supplied record does not specify its smoothness and connectedness conventions for unipotent groups; the intended convention requires confirmation.

Write \(e:\operatorname{Spec}k\to U\) for the identity section. For any \(k\)-scheme \(T\), let \(1_T:T\to U\) denote the composite of its structure morphism with \(e\). All products below are over \(k\).

The proposed rigidity assertion is the following. For every field \(k\), every wound unipotent group \(U\) over \(k\), every closed point \(x\in\mathbb P_k^1\), and every pair \(y_1,y_2\in X(k)\), where
\[
X=\mathbb P_k^1\setminus\{x\},
\]
every \(k\)-morphism \(f:X\times_k X\to U\) satisfying
\[
f\circ i_1=1_X,\qquad f\circ i_2=1_X,
\]
where \(i_1(t)=(y_1,t)\) and \(i_2(t)=(t,y_2)\), satisfies
\[
f=1_{X\times_k X}.
\]
These are equalities of scheme morphisms, rather than merely equalities on rational points. The morphism \(f\) is not assumed to be a group homomorphism. No distinctness condition on \(y_1,y_2\), degree restriction on \(x\), or separability hypothesis on its residue field is supplied.

The record presents this as a conjecture and labels it partially solved. Its supplied dated assessment says that the assertion is proved in characteristic different from two, with structural applications there, while the characteristic-two case and the full all-characteristic assertion remain unresolved. This assessment is not independently verified here; the structural applications are not specified in the record.
\subsection{Short English statement}
The conjecture asks whether a morphism from the product of two copies of a punctured projective line to a wound unipotent group must be constantly the identity whenever it is the identity along both specified coordinate slices. The supplied assessment reports a proof when \(\operatorname{char}(k)\ne 2\), with characteristic two unresolved. The group conventions and dated literature claims require human review.
\subsection{Sources}
\begin{itemize}
\item Supplied local UnsolvedMath registry record, permanent ID \(30004661\), titled \emph{Rigidity of Morphisms to Wound Unipotent Groups}; source problem code \texttt{catalogue entry}, in the Oberwolfach Reports Open Problems collection. The supplied proposed year is 2021; no individual proposer is named.
\item Supplied source citation: \emph{Algebraic Groups (2022)}, Oberwolfach Reports, pp.\ 1087--1148, DOI \texttt{10.4171/owr/2021/20}. \url{https://doi.org/10.4171/owr/2021/20}. The literature-source list calls this \emph{Oberwolfach Report 20/2021}. The discrepancy between the supplied title year and report identifier requires verification.
\item Supplied literature reference: Achet, \emph{Unirational Algebraic Groups}. \url{https://hal.science/hal-02358528/document}.
\item Literature assessment checked at \texttt{2026-08-22}: supplied dated curation, not independently verified. It reports a proof in characteristic different from two and structural applications there, with characteristic two and the full all-characteristic assertion unresolved. The background explicitly requests independent verification of time-sensitive status and cited claims.
\end{itemize}
\end{document}
